% GENERATED FILE. Edit canonical lessons, then run npm run generate:books.
% canonical-source-hash: fnv1a64:5c490f9547b6cfa5
% canonical-lessons: AR-C214-badla, AR-C214-sandal, AR-C214-qumash, AR-C214-qutn, AR-C214-suf

\chapter{Suit, Sandals, Fabric, Cotton, Wool}
\label{ch:ar-suit-sandals-fabric-cotton-wool}
\hlchaptermodality{214}

\begin{chapteropening}
\textbf{By the end of this chapter:} \emph{I can say a suit, sandals, fabric, cotton and wool.}

Complete the last lesson of chapter 214: I can say a suit, sandals, fabric, cotton and wool.
\end{chapteropening}

\section[\texorpdfstring{\ar{بدلة} (badla)}{بدلة (badla)}]{\ar{بدلة} (badla) — a suit}

\label{lesson:AR-C214-badla}



\begin{quote}
Before the new one: say the Arabic for furniture, then the Arabic for laundry.
\end{quote}

\subsection*{You'll want to know: \ar{بدلة}}
\textbf{\ar{بدلة}} — \emph{badla} — \textquotedblleft{}a suit\textquotedblright{}.

\begin{cousinweb}
\textbf{\ar{بدلة}} (\emph{badla}) is built on the native Arabic root \emph{b-d-l}, \textquotedblleft{}to exchange\textquotedblright{} (Lane). \emph{Badla} is a change of clothes, a suit.
\end{cousinweb}

\begin{grammarlens}[title={What you've built}]
One more word for this chapter.
\end{grammarlens}

\subsection*{Your turn}
\begin{itemize}
\raggedright
  \item \emph{Say it:} \emph{badla}
  \item \emph{Say it:} \emph{badla}, once more
  \item \emph{Say it:} say \emph{badla}
  \item \emph{Recall:} say the Arabic for furniture, then the Arabic for laundry, then say \emph{badla} again
\end{itemize}

\begin{tcolorbox}[breakable,colback=teal!4,colframe=teal!35!black,title={Before you move on}]
What does \ar{بدلة} mean? (\textquotedblleft{}A suit\textquotedblright{}.) Say it once more.
\end{tcolorbox}
\section[\texorpdfstring{\ar{صندل} (ṣandal)}{صندل (ṣandal)}]{\ar{صندل} (ṣandal) — sandals}

\label{lesson:AR-C214-sandal}



\begin{quote}
Before the new one: say the Arabic for laundry, then the Arabic for a suit.
\end{quote}

\subsection*{You'll want to know: \ar{صندل}}
\textbf{\ar{صندل}} — \emph{ṣandal} — \textquotedblleft{}sandals\textquotedblright{}.

\begin{cousinweb}
\textbf{\ar{صندل}} (\emph{ṣandal}) is listed by Lane as sandalwood, and his sources also call it a Persian word, arabicized, for a thing like a boot; EtymDB traces the sandalwood word to Middle Persian.
\end{cousinweb}

\begin{grammarlens}[title={What you've built}]
One more word for this chapter.
\end{grammarlens}

\subsection*{Your turn}
\begin{itemize}
\raggedright
  \item \emph{Say it:} \emph{ṣandal}
  \item \emph{Say it:} \emph{ṣandal}, once more
  \item \emph{Say it:} say \emph{ṣandal}
  \item \emph{Recall:} say the Arabic for laundry, then the Arabic for a suit, then say \emph{ṣandal} again
\end{itemize}

\begin{tcolorbox}[breakable,colback=teal!4,colframe=teal!35!black,title={Before you move on}]
What does \ar{صندل} mean? (\textquotedblleft{}Sandals\textquotedblright{}.) Say it once more.
\end{tcolorbox}
\section[\texorpdfstring{\ar{قماش} (qumāsh)}{قماش (qumāsh)}]{\ar{قماش} (qumāsh) — fabric, cloth}

\label{lesson:AR-C214-qumash}



\begin{quote}
Before the new one: say the Arabic for a suit, then the Arabic for sandals.
\end{quote}

\subsection*{You'll want to know: \ar{قماش}}
\textbf{\ar{قماش}} — \emph{qumāsh} — \textquotedblleft{}fabric, cloth\textquotedblright{}.

\begin{cousinweb}
\textbf{\ar{قماش}} (\emph{qumāsh}) is built on the native Arabic root \emph{q-m-sh}, \textquotedblleft{}household stuff\textquotedblright{} (Lane). Lane gives \emph{qammāsh} as a seller of household goods; \emph{qumāsh} is stuff, and fabric is the modern sense.
\end{cousinweb}

\begin{grammarlens}[title={What you've built}]
One more word for this chapter.
\end{grammarlens}

\subsection*{Your turn}
\begin{itemize}
\raggedright
  \item \emph{Say it:} \emph{qumāsh}
  \item \emph{Say it:} \emph{qumāsh}, once more
  \item \emph{Say it:} say \emph{qumāsh}
  \item \emph{Recall:} say the Arabic for a suit, then the Arabic for sandals, then say \emph{qumāsh} again
\end{itemize}

\begin{tcolorbox}[breakable,colback=teal!4,colframe=teal!35!black,title={Before you move on}]
What does \ar{قماش} mean? (\textquotedblleft{}Fabric, cloth\textquotedblright{}.) Say it once more.
\end{tcolorbox}
\section[\texorpdfstring{\ar{قطن} (quṭn)}{قطن (quṭn)}]{\ar{قطن} (quṭn) — cotton}

\label{lesson:AR-C214-qutn}



\begin{quote}
Before the new one: say the Arabic for sandals, then the Arabic for fabric.
\end{quote}

\subsection*{You'll want to know: \ar{قطن}}
\textbf{\ar{قطن}} — \emph{quṭn} — \textquotedblleft{}cotton\textquotedblright{}.

\begin{cousinweb}
\textbf{\ar{قطن}} (\emph{quṭn}) is built on the native Arabic root \emph{q-ṭ-n}, \textquotedblleft{}cotton\textquotedblright{} (Lane). English \emph{cotton} comes from \emph{quṭn} through French.
\end{cousinweb}

\begin{grammarlens}[title={What you've built}]
One more word for this chapter.
\end{grammarlens}

\subsection*{Your turn}
\begin{itemize}
\raggedright
  \item \emph{Say it:} \emph{quṭn}
  \item \emph{Say it:} \emph{quṭn}, once more
  \item \emph{Say it:} say \emph{quṭn}
  \item \emph{Recall:} say the Arabic for sandals, then the Arabic for fabric, then say \emph{quṭn} again
\end{itemize}

\begin{tcolorbox}[breakable,colback=teal!4,colframe=teal!35!black,title={Before you move on}]
What does \ar{قطن} mean? (\textquotedblleft{}Cotton\textquotedblright{}.) Say it once more.
\end{tcolorbox}
\section[\texorpdfstring{\ar{صوف} (ṣūf)}{صوف (ṣūf)}]{\ar{صوف} (ṣūf) — wool}

\label{lesson:AR-C214-suf}



\begin{quote}
Before the new one: say the Arabic for fabric, then the Arabic for cotton.
\end{quote}

\subsection*{You'll want to know: \ar{صوف}}
\textbf{\ar{صوف}} — \emph{ṣūf} — \textquotedblleft{}wool\textquotedblright{}.

\begin{cousinweb}
\textbf{\ar{صوف}} (\emph{ṣūf}) is built on the native Arabic root \emph{ṣ-w-f}, \textquotedblleft{}wool\textquotedblright{} (Lane). An old noun with its own root in Lane.
\end{cousinweb}

\begin{grammarlens}[title={What you've built}]
One more word for this chapter.
\end{grammarlens}

\subsection*{Your turn}
\begin{itemize}
\raggedright
  \item \emph{Say it:} \emph{ṣūf}
  \item \emph{Say it:} \emph{ṣūf}, once more
  \item \emph{Say it:} say \emph{ṣūf}
  \item \emph{Recall:} say the Arabic for fabric, then the Arabic for cotton, then say \emph{ṣūf} again
\end{itemize}

\begin{tcolorbox}[breakable,colback=teal!4,colframe=teal!35!black,title={Before you move on}]
What does \ar{صوف} mean? (\textquotedblleft{}Wool\textquotedblright{}.) Say it once more.
\end{tcolorbox}
